% ============================================================
%  ECON306: W09D2 — Equity: Fair Innings, Distributional Value
%          Judgements, Social Preferences; League Tables;
%          Dementia and End-of-Life Challenges
%  Date: Thursday 17 September 2026
%  SOURCE: COMPRESSION of old w09_week9 (fair innings, ageing)
%    (lines 3105–3165) + old w10_week10 (equity sections)
%    (lines 3178–3380). Best frames per Course Map W09D2 LOs.
%  DROPPED from old deck: Week 10 Overview frame (wrong week),
%    Rawlsianism/Libertarianism detail (combined into one frame here),
%    "Ageing and equity: specific tensions" partial duplication with
%    W09D1 fair innings content. All dropped frames logged in session log.
% ============================================================

\documentclass[aspectratio=169, 12pt]{beamer}

% ============================================================
%  ECON306: Shared Preamble — The Economics of Health and Education
%  University of Otago — Semester 2, 2026
%
%  SHARED BY ALL ECON306 DAY DECKS.
%  Input this file immediately after \documentclass in each deck:
%    % ============================================================
%  ECON306: Shared Preamble — The Economics of Health and Education
%  University of Otago — Semester 2, 2026
%
%  SHARED BY ALL ECON306 DAY DECKS.
%  Input this file immediately after \documentclass in each deck:
%    % ============================================================
%  ECON306: Shared Preamble — The Economics of Health and Education
%  University of Otago — Semester 2, 2026
%
%  SHARED BY ALL ECON306 DAY DECKS.
%  Input this file immediately after \documentclass in each deck:
%    \input{../Preambles/econ306_preamble.tex}
%  Do NOT include \documentclass or per-deck metadata here.
% ============================================================

% ─── Theme & Colours ────────────────────────────────────────
\usetheme{default}
\usecolortheme{default}
\usefonttheme{professionalfonts}
\definecolor{OtagoBlue}{HTML}{003057}
\definecolor{OtagoGold}{HTML}{A8923A}
\definecolor{TextGrey}{HTML}{3A3A3A}
\definecolor{AccentBlue}{HTML}{2B6CB0}
\definecolor{LightBg}{HTML}{F7F9FC}

\setbeamercolor{frametitle}{fg=OtagoBlue, bg=white}
\setbeamercolor{title}{fg=OtagoBlue}
\setbeamercolor{subtitle}{fg=OtagoGold}
\setbeamercolor{author}{fg=TextGrey}
\setbeamercolor{date}{fg=TextGrey}
\setbeamercolor{institute}{fg=TextGrey}
\setbeamercolor{normal text}{fg=TextGrey, bg=white}
\setbeamercolor{item}{fg=OtagoBlue}
\setbeamercolor{subitem}{fg=AccentBlue}
\setbeamercolor{block title}{bg=OtagoBlue!12, fg=OtagoBlue}
\setbeamercolor{block body}{bg=OtagoBlue!5, fg=TextGrey}
\setbeamercolor{alertblock title}{bg=OtagoGold!20, fg=OtagoBlue}
\setbeamercolor{alertblock body}{bg=OtagoGold!8, fg=TextGrey}
\setbeamercolor{alerted text}{fg=OtagoGold!80!black}
\setbeamercolor{structure}{fg=OtagoBlue}

\setbeamertemplate{frametitle}{%
	\vspace{0.35em}%
	\textbf{\insertframetitle}\par%
	\vspace{0.05em}%
	\textcolor{OtagoGold}{\rule{\textwidth}{0.6pt}}}

\setbeamertemplate{navigation symbols}{}
\setbeamertemplate{footline}{%
	\leavevmode\hbox{%
		\begin{beamercolorbox}[wd=.40\paperwidth,ht=2.5ex,dp=1.5ex,leftskip=.3cm]{author in head/foot}%
			\textcolor{TextGrey}{\footnotesize ECON306 \textbullet\ University of Otago}%
		\end{beamercolorbox}%
		\begin{beamercolorbox}[wd=.47\paperwidth,ht=2.5ex,dp=1.5ex,center]{title in head/foot}%
			\textcolor{TextGrey}{\footnotesize \insertshorttitle}%
		\end{beamercolorbox}%
		\begin{beamercolorbox}[wd=.13\paperwidth,ht=2.5ex,dp=1.5ex,rightskip=.3cm,right]{date in head/foot}%
			\textcolor{TextGrey}{\footnotesize \insertframenumber\,/\,\inserttotalframenumber}%
	\end{beamercolorbox}}\vskip0pt}

\setbeamertemplate{itemize item}{\textcolor{OtagoGold}{\small$\blacktriangleright$}}
\setbeamertemplate{itemize subitem}{\textcolor{AccentBlue}{\small$\triangleright$}}
\setbeamertemplate{enumerate item}{\textcolor{OtagoBlue}{\small\bfseries\insertenumlabel.}}

\setbeamertemplate{title page}{%
	\begin{tikzpicture}[remember picture,overlay]
		\fill[OtagoBlue] (current page.north west)
		rectangle ([yshift=-0.38\paperheight]current page.north east);
	\end{tikzpicture}
	\vspace{0.06\paperheight}
	\begin{beamercolorbox}[wd=\textwidth, sep=0pt]{title}
		{\color{white}\Large\textbf{\inserttitle}}\\[0.4em]
		{\color{OtagoGold}\large\insertsubtitle}
	\end{beamercolorbox}
	\vspace{1.4em}
	\begin{beamercolorbox}[wd=\textwidth, sep=0pt]{author}
		{\color{TextGrey}\normalsize\insertauthor}\\[0.2em]
		{\color{TextGrey}\small\insertinstitute}\\[0.2em]
		{\color{TextGrey}\small\insertdate}
\end{beamercolorbox}}

% ─── Packages ──────────────────────────────────────────────
\usepackage{tikz}
\usepackage{booktabs}
\usepackage{array}
\usepackage{graphicx}
\usepackage{hyperref}
\usepackage{amsmath}
\usepackage{amssymb}
\usepackage{colortbl}
\usepackage{multirow}
\usepackage{xcolor}
\hypersetup{colorlinks=true, linkcolor=AccentBlue, urlcolor=AccentBlue,
	citecolor=AccentBlue, pdfauthor={Austin Henderson},
	pdftitle={ECON306 S2 2026}}

% ─── Images folder ─────────────────────────────────────────
%  Place all slide images in the 'images/' subfolder next to each deck.
%  Usage: \includegraphics[width=0.85\textwidth]{figure_name}
\graphicspath{{images/}}

% ─── Reusable macros ───────────────────────────────────────
% Recap slide at start of each lecture
\newcommand{\recapslide}[1]{%
	\begin{frame}{Recap: What Did We Cover Last Time?}
		\begin{block}{Key points from last lecture}
			\begin{itemize}#1\end{itemize}
		\end{block}
\end{frame}}

% Plan slide at start of each lecture
\newcommand{\planslide}[2]{%
	\begin{frame}{Today's Plan — #1}
		\begin{columns}
			\column{0.6\textwidth}
			\begin{itemize}#2\end{itemize}
			\column{0.38\textwidth}
			\begin{alertblock}{Reading}
				\footnotesize See Brightspace (Aoroa) for required readings before this lecture.
			\end{alertblock}
		\end{columns}
\end{frame}}

% Summary slide at end of each lecture
\newcommand{\summaryside}[1]{%
	\begin{frame}{Lecture Summary}
		\begin{exampleblock}{Key takeaways}
			\begin{itemize}#1\end{itemize}
		\end{exampleblock}
		\vspace{0.5em}
		{\footnotesize\textcolor{OtagoGold}{\textbf{Brightspace (Aoroa):}} slides, readings, and any announcements posted there.}
\end{frame}}

% NZ spotlight box
\newcommand{\nzbox}[2]{%
	\begin{alertblock}{\includegraphics[height=0.8em]{nz_flag_icon}\hspace{0.3em} NZ Spotlight: #1}
		#2
\end{alertblock}}

% Tutorial callout
\newcommand{\tutorialbox}[1]{%
	\begin{block}{\textbf{Tutorial This Week}}
		#1
\end{block}}

% ─── Section dividers ──────────────────────────────────────
\AtBeginSection[]{
	\begin{frame}[plain]
		\begin{tikzpicture}[remember picture,overlay]
			\fill[OtagoBlue!95] (current page.north west)
			rectangle (current page.south east);
			\node[anchor=center, text=white, font=\Large\bfseries,
			text width=0.9\paperwidth, align=center]
			at (current page.center) {\insertsectionhead};
			\node[anchor=south, text=OtagoGold, font=\normalsize,
			text width=0.9\paperwidth, align=center]
			at ([yshift=2.5cm]current page.south) {ECON306 — University of Otago};
		\end{tikzpicture}
\end{frame}}

%  Do NOT include \documentclass or per-deck metadata here.
% ============================================================

% ─── Theme & Colours ────────────────────────────────────────
\usetheme{default}
\usecolortheme{default}
\usefonttheme{professionalfonts}
\definecolor{OtagoBlue}{HTML}{003057}
\definecolor{OtagoGold}{HTML}{A8923A}
\definecolor{TextGrey}{HTML}{3A3A3A}
\definecolor{AccentBlue}{HTML}{2B6CB0}
\definecolor{LightBg}{HTML}{F7F9FC}

\setbeamercolor{frametitle}{fg=OtagoBlue, bg=white}
\setbeamercolor{title}{fg=OtagoBlue}
\setbeamercolor{subtitle}{fg=OtagoGold}
\setbeamercolor{author}{fg=TextGrey}
\setbeamercolor{date}{fg=TextGrey}
\setbeamercolor{institute}{fg=TextGrey}
\setbeamercolor{normal text}{fg=TextGrey, bg=white}
\setbeamercolor{item}{fg=OtagoBlue}
\setbeamercolor{subitem}{fg=AccentBlue}
\setbeamercolor{block title}{bg=OtagoBlue!12, fg=OtagoBlue}
\setbeamercolor{block body}{bg=OtagoBlue!5, fg=TextGrey}
\setbeamercolor{alertblock title}{bg=OtagoGold!20, fg=OtagoBlue}
\setbeamercolor{alertblock body}{bg=OtagoGold!8, fg=TextGrey}
\setbeamercolor{alerted text}{fg=OtagoGold!80!black}
\setbeamercolor{structure}{fg=OtagoBlue}

\setbeamertemplate{frametitle}{%
	\vspace{0.35em}%
	\textbf{\insertframetitle}\par%
	\vspace{0.05em}%
	\textcolor{OtagoGold}{\rule{\textwidth}{0.6pt}}}

\setbeamertemplate{navigation symbols}{}
\setbeamertemplate{footline}{%
	\leavevmode\hbox{%
		\begin{beamercolorbox}[wd=.40\paperwidth,ht=2.5ex,dp=1.5ex,leftskip=.3cm]{author in head/foot}%
			\textcolor{TextGrey}{\footnotesize ECON306 \textbullet\ University of Otago}%
		\end{beamercolorbox}%
		\begin{beamercolorbox}[wd=.47\paperwidth,ht=2.5ex,dp=1.5ex,center]{title in head/foot}%
			\textcolor{TextGrey}{\footnotesize \insertshorttitle}%
		\end{beamercolorbox}%
		\begin{beamercolorbox}[wd=.13\paperwidth,ht=2.5ex,dp=1.5ex,rightskip=.3cm,right]{date in head/foot}%
			\textcolor{TextGrey}{\footnotesize \insertframenumber\,/\,\inserttotalframenumber}%
	\end{beamercolorbox}}\vskip0pt}

\setbeamertemplate{itemize item}{\textcolor{OtagoGold}{\small$\blacktriangleright$}}
\setbeamertemplate{itemize subitem}{\textcolor{AccentBlue}{\small$\triangleright$}}
\setbeamertemplate{enumerate item}{\textcolor{OtagoBlue}{\small\bfseries\insertenumlabel.}}

\setbeamertemplate{title page}{%
	\begin{tikzpicture}[remember picture,overlay]
		\fill[OtagoBlue] (current page.north west)
		rectangle ([yshift=-0.38\paperheight]current page.north east);
	\end{tikzpicture}
	\vspace{0.06\paperheight}
	\begin{beamercolorbox}[wd=\textwidth, sep=0pt]{title}
		{\color{white}\Large\textbf{\inserttitle}}\\[0.4em]
		{\color{OtagoGold}\large\insertsubtitle}
	\end{beamercolorbox}
	\vspace{1.4em}
	\begin{beamercolorbox}[wd=\textwidth, sep=0pt]{author}
		{\color{TextGrey}\normalsize\insertauthor}\\[0.2em]
		{\color{TextGrey}\small\insertinstitute}\\[0.2em]
		{\color{TextGrey}\small\insertdate}
\end{beamercolorbox}}

% ─── Packages ──────────────────────────────────────────────
\usepackage{tikz}
\usepackage{booktabs}
\usepackage{array}
\usepackage{graphicx}
\usepackage{hyperref}
\usepackage{amsmath}
\usepackage{amssymb}
\usepackage{colortbl}
\usepackage{multirow}
\usepackage{xcolor}
\hypersetup{colorlinks=true, linkcolor=AccentBlue, urlcolor=AccentBlue,
	citecolor=AccentBlue, pdfauthor={Austin Henderson},
	pdftitle={ECON306 S2 2026}}

% ─── Images folder ─────────────────────────────────────────
%  Place all slide images in the 'images/' subfolder next to each deck.
%  Usage: \includegraphics[width=0.85\textwidth]{figure_name}
\graphicspath{{images/}}

% ─── Reusable macros ───────────────────────────────────────
% Recap slide at start of each lecture
\newcommand{\recapslide}[1]{%
	\begin{frame}{Recap: What Did We Cover Last Time?}
		\begin{block}{Key points from last lecture}
			\begin{itemize}#1\end{itemize}
		\end{block}
\end{frame}}

% Plan slide at start of each lecture
\newcommand{\planslide}[2]{%
	\begin{frame}{Today's Plan — #1}
		\begin{columns}
			\column{0.6\textwidth}
			\begin{itemize}#2\end{itemize}
			\column{0.38\textwidth}
			\begin{alertblock}{Reading}
				\footnotesize See Brightspace (Aoroa) for required readings before this lecture.
			\end{alertblock}
		\end{columns}
\end{frame}}

% Summary slide at end of each lecture
\newcommand{\summaryside}[1]{%
	\begin{frame}{Lecture Summary}
		\begin{exampleblock}{Key takeaways}
			\begin{itemize}#1\end{itemize}
		\end{exampleblock}
		\vspace{0.5em}
		{\footnotesize\textcolor{OtagoGold}{\textbf{Brightspace (Aoroa):}} slides, readings, and any announcements posted there.}
\end{frame}}

% NZ spotlight box
\newcommand{\nzbox}[2]{%
	\begin{alertblock}{\includegraphics[height=0.8em]{nz_flag_icon}\hspace{0.3em} NZ Spotlight: #1}
		#2
\end{alertblock}}

% Tutorial callout
\newcommand{\tutorialbox}[1]{%
	\begin{block}{\textbf{Tutorial This Week}}
		#1
\end{block}}

% ─── Section dividers ──────────────────────────────────────
\AtBeginSection[]{
	\begin{frame}[plain]
		\begin{tikzpicture}[remember picture,overlay]
			\fill[OtagoBlue!95] (current page.north west)
			rectangle (current page.south east);
			\node[anchor=center, text=white, font=\Large\bfseries,
			text width=0.9\paperwidth, align=center]
			at (current page.center) {\insertsectionhead};
			\node[anchor=south, text=OtagoGold, font=\normalsize,
			text width=0.9\paperwidth, align=center]
			at ([yshift=2.5cm]current page.south) {ECON306 — University of Otago};
		\end{tikzpicture}
\end{frame}}

%  Do NOT include \documentclass or per-deck metadata here.
% ============================================================

% ─── Theme & Colours ────────────────────────────────────────
\usetheme{default}
\usecolortheme{default}
\usefonttheme{professionalfonts}
\definecolor{OtagoBlue}{HTML}{003057}
\definecolor{OtagoGold}{HTML}{A8923A}
\definecolor{TextGrey}{HTML}{3A3A3A}
\definecolor{AccentBlue}{HTML}{2B6CB0}
\definecolor{LightBg}{HTML}{F7F9FC}

\setbeamercolor{frametitle}{fg=OtagoBlue, bg=white}
\setbeamercolor{title}{fg=OtagoBlue}
\setbeamercolor{subtitle}{fg=OtagoGold}
\setbeamercolor{author}{fg=TextGrey}
\setbeamercolor{date}{fg=TextGrey}
\setbeamercolor{institute}{fg=TextGrey}
\setbeamercolor{normal text}{fg=TextGrey, bg=white}
\setbeamercolor{item}{fg=OtagoBlue}
\setbeamercolor{subitem}{fg=AccentBlue}
\setbeamercolor{block title}{bg=OtagoBlue!12, fg=OtagoBlue}
\setbeamercolor{block body}{bg=OtagoBlue!5, fg=TextGrey}
\setbeamercolor{alertblock title}{bg=OtagoGold!20, fg=OtagoBlue}
\setbeamercolor{alertblock body}{bg=OtagoGold!8, fg=TextGrey}
\setbeamercolor{alerted text}{fg=OtagoGold!80!black}
\setbeamercolor{structure}{fg=OtagoBlue}

\setbeamertemplate{frametitle}{%
	\vspace{0.35em}%
	\textbf{\insertframetitle}\par%
	\vspace{0.05em}%
	\textcolor{OtagoGold}{\rule{\textwidth}{0.6pt}}}

\setbeamertemplate{navigation symbols}{}
\setbeamertemplate{footline}{%
	\leavevmode\hbox{%
		\begin{beamercolorbox}[wd=.40\paperwidth,ht=2.5ex,dp=1.5ex,leftskip=.3cm]{author in head/foot}%
			\textcolor{TextGrey}{\footnotesize ECON306 \textbullet\ University of Otago}%
		\end{beamercolorbox}%
		\begin{beamercolorbox}[wd=.47\paperwidth,ht=2.5ex,dp=1.5ex,center]{title in head/foot}%
			\textcolor{TextGrey}{\footnotesize \insertshorttitle}%
		\end{beamercolorbox}%
		\begin{beamercolorbox}[wd=.13\paperwidth,ht=2.5ex,dp=1.5ex,rightskip=.3cm,right]{date in head/foot}%
			\textcolor{TextGrey}{\footnotesize \insertframenumber\,/\,\inserttotalframenumber}%
	\end{beamercolorbox}}\vskip0pt}

\setbeamertemplate{itemize item}{\textcolor{OtagoGold}{\small$\blacktriangleright$}}
\setbeamertemplate{itemize subitem}{\textcolor{AccentBlue}{\small$\triangleright$}}
\setbeamertemplate{enumerate item}{\textcolor{OtagoBlue}{\small\bfseries\insertenumlabel.}}

\setbeamertemplate{title page}{%
	\begin{tikzpicture}[remember picture,overlay]
		\fill[OtagoBlue] (current page.north west)
		rectangle ([yshift=-0.38\paperheight]current page.north east);
	\end{tikzpicture}
	\vspace{0.06\paperheight}
	\begin{beamercolorbox}[wd=\textwidth, sep=0pt]{title}
		{\color{white}\Large\textbf{\inserttitle}}\\[0.4em]
		{\color{OtagoGold}\large\insertsubtitle}
	\end{beamercolorbox}
	\vspace{1.4em}
	\begin{beamercolorbox}[wd=\textwidth, sep=0pt]{author}
		{\color{TextGrey}\normalsize\insertauthor}\\[0.2em]
		{\color{TextGrey}\small\insertinstitute}\\[0.2em]
		{\color{TextGrey}\small\insertdate}
\end{beamercolorbox}}

% ─── Packages ──────────────────────────────────────────────
\usepackage{tikz}
\usepackage{booktabs}
\usepackage{array}
\usepackage{graphicx}
\usepackage{hyperref}
\usepackage{amsmath}
\usepackage{amssymb}
\usepackage{colortbl}
\usepackage{multirow}
\usepackage{xcolor}
\hypersetup{colorlinks=true, linkcolor=AccentBlue, urlcolor=AccentBlue,
	citecolor=AccentBlue, pdfauthor={Austin Henderson},
	pdftitle={ECON306 S2 2026}}

% ─── Images folder ─────────────────────────────────────────
%  Place all slide images in the 'images/' subfolder next to each deck.
%  Usage: \includegraphics[width=0.85\textwidth]{figure_name}
\graphicspath{{images/}}

% ─── Reusable macros ───────────────────────────────────────
% Recap slide at start of each lecture
\newcommand{\recapslide}[1]{%
	\begin{frame}{Recap: What Did We Cover Last Time?}
		\begin{block}{Key points from last lecture}
			\begin{itemize}#1\end{itemize}
		\end{block}
\end{frame}}

% Plan slide at start of each lecture
\newcommand{\planslide}[2]{%
	\begin{frame}{Today's Plan — #1}
		\begin{columns}
			\column{0.6\textwidth}
			\begin{itemize}#2\end{itemize}
			\column{0.38\textwidth}
			\begin{alertblock}{Reading}
				\footnotesize See Brightspace (Aoroa) for required readings before this lecture.
			\end{alertblock}
		\end{columns}
\end{frame}}

% Summary slide at end of each lecture
\newcommand{\summaryside}[1]{%
	\begin{frame}{Lecture Summary}
		\begin{exampleblock}{Key takeaways}
			\begin{itemize}#1\end{itemize}
		\end{exampleblock}
		\vspace{0.5em}
		{\footnotesize\textcolor{OtagoGold}{\textbf{Brightspace (Aoroa):}} slides, readings, and any announcements posted there.}
\end{frame}}

% NZ spotlight box
\newcommand{\nzbox}[2]{%
	\begin{alertblock}{\includegraphics[height=0.8em]{nz_flag_icon}\hspace{0.3em} NZ Spotlight: #1}
		#2
\end{alertblock}}

% Tutorial callout
\newcommand{\tutorialbox}[1]{%
	\begin{block}{\textbf{Tutorial This Week}}
		#1
\end{block}}

% ─── Section dividers ──────────────────────────────────────
\AtBeginSection[]{
	\begin{frame}[plain]
		\begin{tikzpicture}[remember picture,overlay]
			\fill[OtagoBlue!95] (current page.north west)
			rectangle (current page.south east);
			\node[anchor=center, text=white, font=\Large\bfseries,
			text width=0.9\paperwidth, align=center]
			at (current page.center) {\insertsectionhead};
			\node[anchor=south, text=OtagoGold, font=\normalsize,
			text width=0.9\paperwidth, align=center]
			at ([yshift=2.5cm]current page.south) {ECON306 — University of Otago};
		\end{tikzpicture}
\end{frame}}


\title[ECON306: W09D2 — Equity and Fair Innings]{ECON306\\The Economics of Health and Education}
\subtitle{W09D2: Equity, Fair Innings, Distributional Value Judgements; League Tables; End-of-Life}
\author{Austin Henderson}
\institute{Department of Economics\\University of Otago}
\date{Thursday 17 September 2026}

\begin{document}

\begin{frame}[plain]
  \titlepage
\end{frame}

\recapslide{
  \item QALY = life years $\times$ quality weight; ICER = $\Delta$cost / $\Delta$QALY
  \item EQ-5D-5L instrument; NZ value set (Devlin, Hansen et al., Otago)
  \item Fund if ICER $<$ threshold; NZ PHARMAC threshold $\approx$\$45--50K/QALY
}

\planslide{Equity and Fairness in Health Allocation}{
  \item QALYs favour younger patients: the fair innings argument
  \item Philosophical frameworks: utilitarianism, Rawlsianism, libertarianism
  \item Common equity criteria: horizontal, vertical, rule of rescue
  \item League tables: strengths and limits
  \item Dementia and end-of-life care challenges for CUA
}

\begin{frame}{QALYs, Ageing, and the ``Fair Innings'' Argument}
  \begin{columns}
    \column{0.54\textwidth}
      \footnotesize
      \begin{itemize}
        \item QALYs are gained over remaining lifetime, so \textbf{younger patients gain more QALYs} from the same treatment
        \item Pure QALY maximisation tends to \textbf{favour treating younger patients}
        \item The \textbf{``fair innings'' argument} (Harris, 1985; Williams, 1997): everyone deserves a ``fair innings'' of healthy life --- prioritising the young may be fair because the old have already had theirs
        \item Counterargument: treating the elderly has intrinsic value regardless of remaining life expectancy
      \end{itemize}
    \column{0.42\textwidth}
      \begin{alertblock}{NZ equity implications}
        \footnotesize M\=aori and Pacific peoples have \textbf{lower life expectancy at 65} than European NZers --- ``fair innings'' logic further disadvantages groups who have had fewer healthy years.\\[0.15em]
        PHARMAC considers equity weights alongside ICERs to address this.
      \end{alertblock}
  \end{columns}
\end{frame}

\begin{frame}{Is QALY Maximisation Enough?}
  \begin{columns}
    \column{0.48\textwidth}
      \begin{block}{The case for QALY maximisation}
        \small Funding in order of cost-per-QALY is \textbf{Pareto efficient}: no other allocation produces more aggregate health.
      \end{block}
    \column{0.48\textwidth}
      \begin{alertblock}{But society also cares about\ldots}
        \footnotesize
        \begin{itemize}
          \item \textbf{Who} benefits? Distribution matters
          \item \textbf{Severity}: worse-off deserve priority
          \item \textbf{Rule of rescue}: save identifiable individuals
          \item \textbf{Fairness}: process and opportunity
          \item \textbf{Innovation}: incentivise R\&D?
          \item \textbf{Carer burden}: unpaid caregiver impacts
        \end{itemize}
      \end{alertblock}
  \end{columns}
\end{frame}

\begin{frame}{Three Philosophical Frameworks}
  \begin{columns}
    \column{0.32\textwidth}
      \begin{block}{Utilitarianism}
        \footnotesize Maximise \textbf{aggregate welfare}. In health: maximise total QALYs.\\[0.15em]
        \textit{Critique}: ignores distribution; inequalities justified if aggregate welfare is high enough
      \end{block}
    \column{0.32\textwidth}
      \begin{alertblock}{Rawlsianism}
        \footnotesize Maximise welfare of the \textbf{worst-off} (maximin). In health: prioritise the sickest.\\[0.15em]
        \textit{Critique}: can justify enormous spending for tiny gains to the worst-off
      \end{alertblock}
    \column{0.32\textwidth}
      \begin{block}{Libertarianism}
        \footnotesize Respect \textbf{individual rights} and voluntary exchange. In health: allow markets, respect patient preferences.\\[0.15em]
        \textit{Critique}: fails those who cannot afford care; ignores externalities
      \end{block}
  \end{columns}
  \vspace{0.15em}
  \begin{exampleblock}{NZ's pragmatic approach}
    \footnotesize PHARMAC applies a mix of utilitarian (QALY max) and Rawlsian (equity weighting for M\=aori and Pacific peoples). No pure framework governs decisions.
  \end{exampleblock}
\end{frame}

\begin{frame}{Common Distributional Value Judgements}
  \begin{enumerate}
    \item \textbf{Horizontal equity}: people with equal needs should receive equal treatment,
          regardless of income, ethnicity, or location
    \item \textbf{Vertical equity}: people with greater needs should receive more care
          (i.e.\ treat the sickest first)
    \item \textbf{Fair innings}: everyone deserves a chance at a full, healthy life ---
          prioritise those who have had fewer healthy years (often the young, or those with
          lower life expectancy)
    \item \textbf{Rule of rescue}: powerful moral intuition to do everything possible for
          identified individuals facing death, even at disproportionate cost
    \item \textbf{Severity weighting}: interventions for more severe conditions receive a
          higher effective threshold (``severity multiplier'', as in Norway)
  \end{enumerate}
\end{frame}

\begin{frame}{Economics Models of Social Preferences}
  \begin{columns}
    \column{0.55\textwidth}
      \begin{block}{Social welfare function (SWF)}
        \small $W = W(U_1, U_2, \ldots, U_n)$. Different forms embed different ethical views:
        \begin{itemize}
          \item \textbf{Utilitarian}: $W = \sum_i U_i$ (maximise sum)
          \item \textbf{Rawlsian}: $W = \min_i U_i$ (maximise minimum)
          \item \textbf{Inequality-averse}: concave $W$ (more weight to worse-off)
        \end{itemize}
      \end{block}
    \column{0.42\textwidth}
      \small
      Distributional weights $w_i$ applied to QALYs:
      \[ W = \sum_i w_i \cdot Q_i \]
      If $w_i > 1$ for disadvantaged groups, their QALYs get extra weight.\\[0.2em]
      \textbf{Challenge}: determine $w_i$ via public preference surveys (vignette studies, discrete choice experiments).
  \end{columns}
\end{frame}

\begin{frame}{The League Table: Strengths and Limitations}
  \begin{itemize}
    \item List all candidate interventions, ranked by \textbf{cost per QALY} (ICER)
    \item Fund from the top down until the budget is exhausted
    \item \textbf{Appeal}: transparent, consistent, maximises aggregate health
    \item \textbf{Problems with league tables}:
    \begin{itemize}
      \item Different studies use different methods (comparators, perspectives, time horizons)
      \item Assumes all QALYs are equal (no distributional weights)
      \item Ignores budget impact (a \$5M drug with a low ICER exhausts the budget)
      \item Does not account for uncertainty in ICER estimates
      \item Ignores non-QALY factors: equity, innovation, unmet need
    \end{itemize}
    \item \textbf{PHARMAC does not use a formal league table} --- it weighs multiple criteria
          (including equity, clinical benefit, budget impact) in a committee process
  \end{itemize}
\end{frame}

\begin{frame}{Dementia and End-of-Life Care: Challenges for CUA}
  \small
  \begin{columns}
    \column{0.55\textwidth}
      \textbf{Dementia and the QALY framework:}
      \begin{itemize}
        \item Cognitive impairment makes self-reported EQ-5D unreliable (proxy responses needed)
        \item Proxy utility weights may not reflect the patient's subjective experience
        \item Dementia care costs are enormous but QALY gains from care are modest
      \end{itemize}
    \column{0.42\textwidth}
      \textbf{End-of-life care:} NZ spends disproportionately in the last months of life, with limited QALY gains per dollar.\\[0.25em]
      \textbf{NICE (2009)}: end-of-life premium (£50K vs.\ £30K threshold).\\[0.25em]
      NZ PHARMAC has no formal end-of-life premium --- ongoing policy debate. As NZ ages, the QALY framework needs adaptation.
  \end{columns}
\end{frame}

\summaryside{
  \item QALYs favour younger patients; fair innings argument partially justifies age-weighting but disadvantages groups with shorter life expectancy (M\=aori, Pacific)
  \item Three frameworks: utilitarianism (max aggregate QALYs), Rawlsianism (maximin), libertarianism (rights and markets)
  \item Distributional weights ($w_i > 1$ for disadvantaged groups) and social welfare functions formalise equity concerns
  \item League tables are transparent but ignore equity, budget impact, and uncertainty; PHARMAC uses a multi-criteria committee process
}

\end{document}
